% TOP_PROBLEM: {"id": 30001846, "title": "The Abelian Rank-One Stark Conjecture", "category_id": 1, "category": {"id": 1, "name": "number_theory", "display_name": "Number Theory", "description": "Properties of integers, prime numbers, Diophantine equations.", "slug": "number-theory", "order_index": 1, "created_at": "2026-07-31T15:26:25.670Z"}, "status": "partially_solved"}
% FIRST_POSTED: 2026-09-27
% !TeX program = lualatex
% Compile twice: lualatex open_problems_500.tex
% All references and prose are embedded; no BibTeX database or external inputs.
\documentclass[11pt,a4paper]{article}
\usepackage[margin=24mm,headheight=14pt]{geometry}
\usepackage{fontspec}
\setmainfont{Latin Modern Roman}
\setsansfont{Latin Modern Sans}
\setmonofont{Latin Modern Mono}
\usepackage{amsmath,amssymb,mathtools}
\usepackage{unicode-math}
\setmathfont{Latin Modern Math}
\usepackage{microtype}
\usepackage{enumitem,xurl,needspace}
\usepackage[dvipsnames]{xcolor}
\usepackage{hyperref}
\hypersetup{unicode=true,colorlinks=true,linkcolor=MidnightBlue,urlcolor=MidnightBlue,
 pdftitle={500 Mathematical Problem Statements},pdfauthor={Source-annotated compilation},bookmarksnumbered=true}
\usepackage{bookmark}
\usepackage{tocloft}
\setlength{\cftsecnumwidth}{3em}
\cftsetpnumwidth{2.5em}
\cftsetrmarg{4em}
\usepackage{titlesec}
\titleformat{\section}{\Large\bfseries\raggedright}{\thesection}{0.7em}{}
\usepackage{fancyhdr}
\pagestyle{fancy}\fancyhf{}
\fancyhead[L]{\small\sffamily Mathematical problem statements}
\fancyhead[R]{\small Source edition 22}
\fancyfoot[C]{\thepage}
\setlength{\parindent}{0pt}
\setlength{\parskip}{5pt plus 1pt}
\setlength{\emergencystretch}{3em}
\setcounter{tocdepth}{1}
\setcounter{secnumdepth}{1}
\setlist[itemize]{leftmargin=3em,itemsep=5pt,topsep=4pt,parsep=0pt}
\allowdisplaybreaks
\newcommand{\N}{\mathbb{N}}
\newcommand{\Z}{\mathbb{Z}}
\newcommand{\Q}{\mathbb{Q}}
\newcommand{\R}{\mathbb{R}}
\newcommand{\C}{\mathbb{C}}
\newcommand{\F}{\mathbb{F}}
\newcommand{\A}{\mathbb{A}}
\newcommand{\PP}{\mathbb P}
\newcommand{\E}{\mathbb{E}}
\newcommand{\eps}{\varepsilon}
\newcommand{\dd}{\,\mathrm d}
\newcommand{\sref}[2]{\hyperlink{source:#1:#2}{[S#2]}}
\newcommand{\eref}[2]{\hyperlink{extra:#1:#2}{[E#2]}}
% Turn the next switch off for a shorter reading copy. The quotations preserve
% every exactTarget field, including qualifications omitted from a short summary.
\newif\ifshowcatalogue
\showcataloguetrue
\newcommand{\cataloguescope}[1]{\ifshowcatalogue
 \paragraph{Exact scope in the supplied catalogue.}
 \begin{quote}\small #1\end{quote}\fi}

% Compile twice with: lualatex 30001846.tex
\numberwithin{equation}{section}
\hypersetup{pdftitle={Research attempt -- problem 30001846},pdfauthor={Research notebook; original compilation retained}}
\fancyhead[L]{\small\sffamily Research notebook}
\fancyhead[R]{\small\sffamily Problem 30001846 | 27 September 2026}
\begin{document}
\setcounter{section}{0}
% ================================================================
% problemId: 30001846
% Canonical problem ID: 30001846
% exactTarget SHA-256: 59f04a7ed7ef6089b0ff1799fcdd7ae6ac5d164ee39f961373f713b675ca4a10
\clearpage
\section*{\texorpdfstring{The Abelian Rank-One Stark Conjecture}{The Abelian Rank-One Stark Conjecture}}\label{30001846}
\subsection{Definitions and mathematical statement}
For an abelian extension $K/F$ and $S$ as specified, the partial zeta function $\zeta_{K/F,S}(\sigma,s)$ sums $(N\mathfrak a)^{-s}$ over ideals of $F$ prime to $S$ having Artin symbol $\sigma$, continued to $s=0$. The group $U_{v,S}$ of permitted Stark units, its normalized local absolute values, and the Artin-symbol action convention are those of the cited notes; these conventions affect the displayed identity. If $w\mid v$ and $e=|\mu(K)|$, the requested unit satisfies
\[
 \zeta'_{K/F,S}(\sigma,0)=-e^{-1}\log|u^\sigma|_w
 \quad(\sigma\in G),\qquad K(u^{1/e})/F\text{ abelian}.
\]
The last condition is part of the conjecture, not a consequence to be dropped from a regulator-only reformulation. The source definition of $U_{v,S}$ is retained rather than replacing it with unrestricted units.
\par\smallskip\textit{Formulation and concept references:} \sref{30001846}{1}.
\cataloguescope{Let K/F be an abelian extension of number fields with Galois group G, and let S be a finite set of places of F containing the archimedean and ramified places, with |S| >= 2 and at least one place v of S splitting completely in K. Prove that, for a fixed place w of K above v, there exists a unit u in U\_\{v,S\} such that zeta'\_\{K/F,S\}(sigma, 0) = -(1/e) log |u\textasciicircum{}sigma|\_w for all sigma in G, where e is the number of roots of unity in K, and K(u\textasciicircum{}\{1/e\})/F is abelian.}
\subsection{Short English statement}
Can all the relevant partial-zeta derivatives be realized by logarithms of conjugates of a single permitted unit whose root generates an abelian extension?
% BEGIN RESEARCH ADDITION ATTEMPT-198
\subsection{Research attempt}

\paragraph{Current frontier.}
The formulation is Conjecture 1.1 of the source, including the permitted-unit
condition and the abelian radical extension. Its Section 1.4 separates
archimedean and finite-place cases. Dasgupta--Kakde prove the Brumer--Stark
conjecture away from 2; the subsequent Dasgupta--Kakde--Silliman--Wang
manuscript gives a proof over $\Z$ in the finite-place, totally-real/CM
setting. The latter manuscript's statements were inspected here, not its
entire proof independently audited. Neither statement should be relabelled
as a proof for every archimedean split place in the present question.
\sref{30001846}{1}, \eref{30001846}{1}, \eref{30001846}{2}.
Honnor's work on root-of-unity ambiguity in a Brumer--Stark formula also
underscores that obtaining logarithms and obtaining the specified radical
extension are different tasks. \eref{30001846}{3}.

\paragraph{Attack route.}
Separate three possible losses: existence in a real regulator space,
integrality in the unit lattice, and the commutator obstruction for the
radical extension. A regulator argument controls the first, rational Stark
statements can control a rational lattice, and Kummer theory controls the
last; none automatically supplies the missing integral point. The route
pursued here makes these obstructions independently testable. The Kummer
criterion below is proved under an explicit full-degree hypothesis; it is
not asserted to cover the lower-degree cases without modification.

\paragraph{Attempt.}
\textbf{1. The exact lattice, not unrestricted $S$-units.}
The source defines $U=U_{v,S}$ as follows. If $|S|\ge3$, require
$|u|_{w'}=1$ at every place not above $v$. If $S=\{v,v'\}$, require the
absolute values at places above $v'$ to be equal to one another, and
$|u|_{w'}=1$ outside $S$. The values are the source's product-formula
normalized ones. \sref{30001846}{1}, Section 1.2, printed page 9.
Set
\begin{equation}\label{eq:30001846:lattice}
 \lambda(u)=(\log|u^\sigma|_w)_{\sigma\in G},\qquad
 b=(-e\zeta'_{K/F,S}(\sigma,0))_{\sigma\in G}.
\end{equation}
Then $\ker\lambda=\mu(K)$. Indeed, vanishing of these coordinates gives
vanishing at every place above $v$, because $v$ splits completely.
For $|S|\ge3$ all other logarithms vanish by definition. For $|S|=2$
the product formula forces the common logarithm above $v'$ to vanish as
well. An algebraic number with absolute value one at every place is a
root of unity: it is an algebraic integer unit, and all conjugates have
modulus one, so the bounded-coefficient argument applied to its powers
proves Kronecker's assertion.

More precisely, the monic polynomials whose roots are the conjugates of
$u^m$ have integral coefficients bounded independently of $m$. There are
only finitely many such polynomials and hence finitely many possible
values of $u^m$. Two powers coincide, proving that $u$ is torsion.
Thus the kernel assertion does not use Stark's conjecture.

The full $S$-unit logarithms form a lattice modulo torsion by the
$S$-unit theorem. On the real subspace specified by the preceding zero
and equal-coordinate conditions, projection to the coordinates above
$v$ has a continuous linear inverse onto its image: the omitted common
coordinate, when present, is recovered by the product formula.
Consequently $\lambda(U)$ is discrete in its real span. The logarithmic
part of the target is the exact membership test
\begin{equation}\label{eq:30001846:membership}
                         b\in\lambda(U).
\end{equation}
Showing $b\in\R\lambda(U)$ is much weaker. Even showing
$b\in\Q\lambda(U)$ proves only $db=\lambda(u_0)$ for some integer
$d>0$; taking a $d$th root of $u_0$ may leave $K$. All solutions of
\eqref{eq:30001846:membership}, when any exists, are $u\xi$ with
$\xi\in\mu(K)$. These assertions concern the original unit group, not
an enlarged real or rational group.

\textbf{2. An explicit abelianity certificate.}
Suppose $u\in U$, $\beta^e=u$, and $[K(\beta):K]=e$. Let
$c_\sigma\in\Z$ represent the action of $\sigma$ on $\mu_e$:
$\sigma(\zeta)=\zeta^{c_\sigma}$. The following two conditions are
equivalent to $K(\beta)/F$ being abelian:
\begin{align}
 \sigma(u)&=\gamma_\sigma^e u^{c_\sigma}
       &&\text{for some }\gamma_\sigma\in K^*,\quad\sigma\in G,
       \label{eq:30001846:power}\\
 \sigma(\gamma_\tau)\gamma_\sigma^{c_\tau}
       &=\tau(\gamma_\sigma)\gamma_\tau^{c_\sigma}
       &&\text{for every }\sigma,\tau\in G.
       \label{eq:30001846:commutator}
\end{align}
Here the representatives $c_\sigma$ are fixed before choosing the
$\gamma_\sigma$.

For sufficiency, \eqref{eq:30001846:power} defines an extension
$S_\sigma$ of $\sigma$ by
$S_\sigma(\beta)=\gamma_\sigma\beta^{c_\sigma}$.
Irreducibility of $T^e-u$ and the displayed $e$th-power identity make
this an $F$-automorphism of $K(\beta)$. Such lifts for every $\sigma$,
together with $\beta\mapsto\zeta\beta$, make the extension Galois.
Each lift commutes with this cyclic kernel, since both orders of
composition send $\beta$ to
$\gamma_\sigma\zeta^{c_\sigma}\beta^{c_\sigma}$.
The two compositions $S_\sigma S_\tau$ and $S_\tau S_\sigma$ agree
on $K$ and send $\beta$ respectively to
\[
 \sigma(\gamma_\tau)\gamma_\sigma^{c_\tau}
       \beta^{c_\sigma c_\tau},\qquad
 \tau(\gamma_\sigma)\gamma_\tau^{c_\sigma}
       \beta^{c_\sigma c_\tau}.
\]
Equation \eqref{eq:30001846:commutator} therefore makes all the generators
commute.

Conversely, if the extension is abelian, choose any lift $S_\sigma$.
Commutation with every $\beta\mapsto\zeta\beta$ says that
$S_\sigma(\beta)/\beta^{c_\sigma}$ is fixed by
$\operatorname{Gal}(K(\beta)/K)$, hence is some
$\gamma_\sigma\in K^*$. This gives \eqref{eq:30001846:power}; commutation
of the lifts gives \eqref{eq:30001846:commutator}. Moreover, under
\eqref{eq:30001846:power} the ratio of the two sides of
\eqref{eq:30001846:commutator} belongs to $\mu_e$. Replacing a
$\gamma_\sigma$ by a root-of-unity multiple does not change that
ratio, because $\sigma(\zeta)=\zeta^{c_\sigma}$. The obstruction is
therefore not removed merely by changing the chosen lifts.

\textbf{3. Normalization test.}
For $K=F=\Q$, $S=\{\infty,p\}$, one has $e=2$ and
\[
 \zeta_S(s)=(1-p^{-s})\zeta(s),\qquad
                  \zeta'_S(0)=-\tfrac12\log p.
\]
With $v=\infty$ take $u=p$; with $v=p$ take $u=p^{-1}$, since
$|p^{-1}|_p=p$. Both choices satisfy the source's $U_{v,S}$ condition,
and their square roots generate quadratic, hence abelian, extensions.
The computation checks both the sign and the factor $e$; it does not
produce units over general $F$.

\paragraph{Stress test.}
A central extension of an abelian group need not be abelian: the quaternion
group maps onto the Klein four group with central kernel. Thus the
power congruences alone cannot justify deleting
\eqref{eq:30001846:commutator}. This group example is an obstruction to an
inference, not a counterexample to Stark's statement. Changing $u$ by
$\xi\in\mu(K)$ preserves every logarithm but may change its Kummer
class. One must find a suitable representative, not choose one
arbitrarily. Honnor's result concerns this ambiguity in its stated
finite-place formula, not all instances of \eqref{eq:30001846:membership}.
\eref{30001846}{3}.

The full-degree assumption in the certificate is essential to the proof
as written: it identifies the entire cyclic kernel with $\mu_e$.
No claim about all lower-degree radicals has been smuggled into the
criterion. Most importantly, no argument above proves the lattice
membership \eqref{eq:30001846:membership}, even before abelianity is imposed.

\paragraph{Outcome.}
\textbf{Outcome: Conditional reduction.}
The original logarithmic requirement is isolated as an integral lattice
problem with exactly root-of-unity ambiguity. In the full-degree Kummer
case, equations \eqref{eq:30001846:power}--\eqref{eq:30001846:commutator} give a
proved certificate for the second requirement. A resolution still needs
the exact lattice point, a representative with the required abelianity,
and treatment of the remaining Kummer degrees. The rational leading-term
problem in Section 166 is not a substitute for these integral conditions.

% END RESEARCH ADDITION ATTEMPT-198
\subsection{Sources}
\begin{itemize}\raggedright
\item[\textnormal{[S1]}]\hypertarget{source:30001846:1}{} S. Dasgupta and M. Greenberg, The Rank One Abelian Stark Conjecture, Arizona Winter School 2011 course notes, March 11, 2011.
\url{https://swc-math.github.io/aws/2011/2011DasguptaGreenbergNotes.pdf}.
{\small\textit{Catalogue formulation source.}}
% BEGIN RESEARCH ADDITION REFERENCES-198
\item[\textnormal{[E1]}]\hypertarget{extra:30001846:1}{} Samit Dasgupta and Mahesh Kakde, \emph{On the Brumer--Stark conjecture}, Annals of Mathematics 197 (2023), 289--388; arXiv:2010.00657. Main theorem over $\Z[1/2]$ and its precise totally-real/CM setting.
\url{https://arxiv.org/abs/2010.00657}.
{\small\textit{Research reference; consulted 27 September 2026.}}
\item[\textnormal{[E2]}]\hypertarget{extra:30001846:2}{} Samit Dasgupta, Mahesh Kakde, Jesse Silliman, and Jiuya Wang, \emph{The Brumer--Stark Conjecture over $\Z$}, authors\textquotesingle{} manuscript, 56 pages. Section 1.1, Theorem 1.1 and Section 1.2, Theorem 1.2.
\url{https://sites.math.duke.edu/~dasgupta/papers/Brumer-Stark-Z.pdf}.
{\small\textit{Research reference; statement and scope inspected 27 September 2026; the complete proof was not independently audited here.}}
\item[\textnormal{[E3]}]\hypertarget{extra:30001846:3}{} Matthew H. Honnor, \emph{On the Root of Unity Ambiguity in a Formula for the Brumer--Stark Units}, arXiv:2302.12332v3, Introduction and the stated root-of-unity refinement.
\url{https://arxiv.org/html/2302.12332v3}.
{\small\textit{Research reference; consulted 27 September 2026.}}
% END RESEARCH ADDITION REFERENCES-198
\end{itemize}

\end{document}
